\documentclass[10pt,a4paper]{article}

% ---------- Encoding / fonts ----------
\usepackage{fontspec}
\usepackage{polyglossia}
\setmainlanguage{russian}
\setotherlanguage{english}
\setmainfont{Liberation Sans}
\usepackage{graphicx}

% ---------- Layout ----------
\usepackage[margin=1.25cm,top=0.55cm,bottom=0.55cm]{geometry}
\usepackage{titlesec}
\usepackage{enumitem}
\usepackage{hyperref}
\usepackage{xcolor}
\usepackage{parskip}
\usepackage{microtype}
\sloppy
\setlength{\parskip}{0.05em}

\definecolor{accent}{HTML}{2B4C7E}
\definecolor{darkgray}{HTML}{333333}
\definecolor{midgray}{HTML}{666666}

\hypersetup{
  colorlinks=true,
  urlcolor=accent,
  linkcolor=accent,
  pdftitle={Алексеев Даниил --- DevOps Intern / SRE Intern},
  pdfauthor={Алексеев Даниил Александрович}
}

\pagestyle{empty}
\setlength{\parindent}{0pt}

% ---------- Section styling ----------
\titleformat{\section}
  {\large\bfseries\color{accent}}
  {}{0em}{}
  [\vspace{-0.1em}{\color{accent!40}\titlerule}]
\titlespacing*{\section}{0pt}{0.2em}{0.15em}

\titleformat{\subsection}
  {\bfseries\color{darkgray}}
  {}{0em}{}
\titlespacing*{\subsection}{0pt}{0.25em}{0.05em}

\newcommand{\jobmeta}[1]{{\small\color{midgray}\textit{#1}}\par\vspace{0.05em}}
\newcommand{\stackline}[1]{\par{\small\color{midgray}Стек: #1}\par}
\newcommand{\repoline}[1]{\par{\small\color{midgray}Репозиторий: #1}\par}
% Logos use the same paths as the other CVs. Missing logos do not break compilation.
\newcommand{\universitylogo}[1]{\IfFileExists{#1}{\includegraphics[width=1.2cm]{#1}}{}}

\setlist[itemize]{leftmargin=1.1em, itemsep=0.05em, topsep=0.1em, parsep=0pt}

\begin{document}

% ----- Header -----
\begin{center}
  {\LARGE\bfseries Алексеев Даниил Александрович}\\[0.15em]
  {\normalsize\color{accent} DevOps Intern / SRE Intern}\\[0.3em]
  {\small
    \href{mailto:dalyoksi@yandex.ru}{dalyoksi@yandex.ru}
    \;\textbullet\; +79114202996 \;\textbullet\; Москва
    \;\textbullet\; \href{https://github.com/alyoksi}{github.com/alyoksi}
  }
\end{center}
\vspace{-0.3em}

% ----- About -----
\section*{О себе}
Магистрант МГТУ им. Н.Э. Баумана, ищу стажировку в DevOps/SRE. Интересуюсь компьютерными сетями, автоматизацией развёртывания и отказоустойчивостью распределённых систем. Пятый год использую Linux как настольную ОС: Debian и дистрибутивы на базе Ubuntu.

% ----- Skills ----
\section*{Технический стек}

\textbf{Языки:} Python, C/C++, C\#, Go, Lua

\textbf{Базы данных:} PostgreSQL, Redis, MariaDB, Neo4j

\textbf{DevOps \& CI/CD:} Docker, Docker Compose, GitHub Actions, Coveralls, SonarCloud; Kubernetes, Helm, Ingress

\textbf{Сборка \& Окружение:} Linux, Git, CMake, Make, GDB

\textbf{Сети:} TCP/IP, ARP, ICMP, DNS, DHCP, HTTP(S); VLAN, NAT/PAT, статическая маршрутизация; принципы STP, OSPF и BGP

% ----- Projects -----
\section*{Проекты}

\subsection*{Развёртывание и отказоустойчивость микросервисной системы}
\jobmeta{Учебный проект по курсу «Распределённые системы обработки информации», МГТУ им. Н.Э. Баумана}
\begin{itemize}
  \item Реализовал на Go систему взаимодействующих по HTTP/REST сервисов с единым API Gateway и отдельными логическими базами PostgreSQL.
  \item Контейнеризировал сервисы и настроил локальную сборку и запуск приложения через Docker Compose.
  \item Реализовал Circuit Breaker и fallback-ответы; откат операций при сбоях и очередь повторной обработки запросов.
  \item Настроил в GitHub Actions сборку и тестирование приложения, сборку и публикацию Docker-образов, автоматический деплой сервисов в Kubernetes.
  \item Развернул приложение в managed-кластере Kubernetes; подготовил параметризованный Helm chart для сервисов и настроил публикацию через Ingress.
  \item Проверил взаимодействие сервисов интеграционными тестами и поведение системы при недоступности отдельных компонентов.
\end{itemize}
\stackline{Go, PostgreSQL, Docker, Docker Compose, GitHub Actions, Kubernetes, Helm, Ingress, Postman}
% Add a link to your actual Go implementation here before sending the CV.
% \repoline{\href{https://github.com/USERNAME/REPOSITORY}{github.com/USERNAME/REPOSITORY}}

\subsection*{Тестирующая система для олимпиад по программированию (грантовый проект)}
\jobmeta{Команда: 4 человека}
\begin{itemize}
  \item Разработал REST API на C\# для обработки решений, авторизации и взаимодействия с базой данных.
  \item Настроил контейнеризацию сервисов с помощью Docker и развёртывание системы; проект использовался на учебных сборах (около 30 участников).
\end{itemize}
\stackline{C\#, REST API, Docker, PostgreSQL, React}

% ----- Experience -----
\section*{Опыт работы}

\subsection*{Разработчик --- оплачиваемая учебная практика}
\jobmeta{3 месяца, вуз}
Разработал приложение для расчёта параметров раскроя заготовки по STL/DXF-модели, автоматизировав расчёт, ранее выполнявшийся вручную. Реализовал потоковый парсер STL, алгоритмы вычислительной геометрии и GUI с 3D-визуализацией и кэшированием геометрии.
\stackline{Python, NumPy, Numba, SciPy, PySide6, VTK}
\repoline{\href{https://github.com/alyoksi/stl-metrics}{github.com/alyoksi/stl-metrics}}

\subsection*{Лектор --- Олимпиадные школы МФТИ}
\jobmeta{Сезонные смены: две летние и одна зимняя}
Читал лекции по алгоритмам и структурам данных для школьников и студентов, составлял задачи для тренировок и туров.

% ----- Achievements -----
\section*{Достижения}
\begin{itemize}
  \item Дважды призёр полуфинала ICPC.
  \item Призёр Чемпионата Урала по программированию, дважды призёр Чемпионата RuCode.
\end{itemize}

% ----- Education -----
\section*{Образование}

\begin{minipage}[c]{0.08\textwidth}
    \centering
    \universitylogo{../images/bmstu_logo.jpg}
\end{minipage}%
\begin{minipage}[c]{0.92\textwidth}
    \textbf{Магистратура, Программная инженерия (09.04.04)} --- МГТУ им. Н.Э. Баумана, Москва, 2026--2028 (текущее обучение)\\
    \textit{Специализация:} Программно-алгоритмическое обеспечение систем ИИ
\end{minipage}

\vspace{0.4em}

\begin{minipage}[c]{0.08\textwidth}
    \centering
    \universitylogo{../images/petrsu_logo.png}
\end{minipage}%
\begin{minipage}[c]{0.92\textwidth}
    \textbf{Бакалавриат, Программная инженерия (09.03.04)} --- Петрозаводский государственный университет, выпуск 2026\\
    \textit{Профиль:} Системное и прикладное программное обеспечение
\end{minipage}

\vspace{0.5em}
\textbf{Доп. образование:} Компьютерные сети (видеолекции и онлайн-курс по eNSP).

% ----- Languages -----
\section*{Языки}
Английский --- B2 \quad\textbullet\quad Финский --- A1

\end{document}
